% =================================================
% Краткие ответы для самостоятельной проверки
% Включаются только в ученическую версию практикума.
% =================================================
\newpage
\homeworktitle
  {Краткие ответы}
  {Проверяйте только после самостоятельного выполнения}

\begin{selfcheck}[Как пользоваться ответами]
  Сначала завершите решение и запишите ответ. Затем сравните результат с
  ключом. Если ответы не совпали, вернитесь к последнему переходу, в котором
  вы были уверены. Полные решения остаются у учителя.
\end{selfcheck}

\begin{answerkeybox}[Задание 1]
  а) \(\lg^3 A=(\lg A)^3\); в \(\lg(A^3)\) число \(3\) относится к
  аргументу; в \(\log_3 A\) число \(3\) является основанием.

  б) \(\lg^2 1000=9\), \(\lg(1000^2)=6\).
\end{answerkeybox}

\begin{answerkeybox}[Задание 2]
  Левая часть: \(x\ne1\). Правая часть: \(x>1\).
  \[
    \lg((x-1)^2)=2\lg|x-1|,\qquad x\ne1.
  \]
\end{answerkeybox}

\begin{answerkeybox}[Задание 3]
  \[
    y\in(-\infty,-2]\cup[2,+\infty).
  \]
\end{answerkeybox}

\begin{answerkeybox}[Задание 4]
  а) \(x\in[-1,2)\cup(2,5]\).

  б)
  \[
    x\in[-\sqrt5,-2)\cup(-2,-\sqrt3]
    \cup[\sqrt3,2)\cup(2,\sqrt5].
  \]
\end{answerkeybox}

\begin{answerkeybox}[Задание 5]
  \[
    x\in(-\infty,-3]
    \cup\left[\frac34,1\right)
    \cup\left(1,\frac54\right]
    \cup[5,+\infty).
  \]
\end{answerkeybox}

\begin{answerkeybox}[Задание 6]
  Вариант А:
  \[
    t=\lg^2|x^2-1|=\bigl(\lg|x^2-1|\bigr)^2.
  \]
\end{answerkeybox}

\begin{answerkeybox}[Задание 7]
  \[
    x\in(-\infty,-11]
    \cup\left[-\frac{19}{9},-2\right)
    \cup\left(-2,-\frac{17}{9}\right]
    \cup[7,+\infty).
  \]
\end{answerkeybox}

\begin{answerkeybox}[Задание 8]
  \[
    \begin{aligned}
      x\in{}&(-\infty,-\sqrt5]
      \cup\left[-\frac{\sqrt5}{2},-1\right)
      \cup\left(-1,-\frac{\sqrt3}{2}\right]\\
      &\cup\left[\frac{\sqrt3}{2},1\right)
      \cup\left(1,\frac{\sqrt5}{2}\right]
      \cup[\sqrt5,+\infty).
    \end{aligned}
  \]
\end{answerkeybox}

\begin{answerkeybox}[Задание 9]
  \[
    A\in(-\infty,-b^2]
    \cup[-b^{-2},0)
    \cup(0,b^{-2}]
    \cup[b^2,+\infty).
  \]
\end{answerkeybox}
